\documentclass[10pt,a4paper]{amsart}
\usepackage{amsmath,amssymb,amsthm,graphicx}
\usepackage{verbatim}
\newenvironment{reference}{\comment}{\endcomment}
%\newenvironment{reference}{}{}
\newenvironment{slogan}{\comment}{\endcomment}
\newenvironment{history}{\comment}{\endcomment}

% For commutative diagrams we use Xy-pic
\usepackage[all]{xy}

% We use 2cell for 2-commutative diagrams.
\xyoption{2cell}
% Theorem environments.
%
\theoremstyle{plain}
\newtheorem{theorem}[subsection]{Theorem}
\newtheorem{proposition}[subsection]{Proposition}
\newtheorem{lemma}[subsection]{Lemma}

\theoremstyle{definition}
\newtheorem{definition}[subsection]{Definition}
\newtheorem{example}[subsection]{Example}
\newtheorem{exercise}[subsection]{Exercise}
\newtheorem{situation}[subsection]{Situation}

\theoremstyle{remark}
\newtheorem{remark}[subsection]{Remark}
\newtheorem{remarks}[subsection]{Remarks}

\numberwithin{equation}{subsection}

% Macros
%
\def\lim{\mathop{\mathrm{lim}}\nolimits}
\def\colim{\mathop{\mathrm{colim}}\nolimits}
\def\Spec{\mathop{\mathrm{Spec}}}
\def\Hom{\mathop{\mathrm{Hom}}\nolimits}
\def\Ext{\mathop{\mathrm{Ext}}\nolimits}
\def\SheafHom{\mathop{\mathcal{H}\!\mathit{om}}\nolimits}
\def\SheafExt{\mathop{\mathcal{E}\!\mathit{xt}}\nolimits}
\def\Sch{\mathit{Sch}}
\def\Mor{\mathop{\mathrm{Mor}}\nolimits}
\def\Ob{\mathop{\mathrm{Ob}}\nolimits}
\def\Sh{\mathop{\mathit{Sh}}\nolimits}
\def\NL{\mathop{N\!L}\nolimits}
\def\CH{\mathop{\mathrm{CH}}\nolimits}
\def\proetale{{pro\text{-}\acute{e}tale}}
\def\etale{{\acute{e}tale}}
\def\QCoh{\mathit{QCoh}}
\def\Ker{\mathop{\mathrm{Ker}}}
\def\Im{\mathop{\mathrm{Im}}}
\def\Coker{\mathop{\mathrm{Coker}}}
\def\Coim{\mathop{\mathrm{Coim}}}

% Boxtimes
%
\DeclareMathSymbol{\boxtimes}{\mathbin}{AMSa}{"02}

%
% Macros for moduli stacks/spaces
%
\def\QCohstack{\mathcal{QC}\!\mathit{oh}}
\def\Cohstack{\mathcal{C}\!\mathit{oh}}
\def\Spacesstack{\mathcal{S}\!\mathit{paces}}
\def\Quotfunctor{\mathrm{Quot}}
\def\Hilbfunctor{\mathrm{Hilb}}
\def\Curvesstack{\mathcal{C}\!\mathit{urves}}
\def\Polarizedstack{\mathcal{P}\!\mathit{olarized}}
\def\Complexesstack{\mathcal{C}\!\mathit{omplexes}}
% \Pic is the operator that assigns to X its picard group, usage \Pic(X)
% \Picardstack_{X/B} denotes the Picard stack of X over B
% \Picardfunctor_{X/B} denotes the Picard functor of X over B
\def\Pic{\mathop{\mathrm{Pic}}\nolimits}
\def\Picardstack{\mathcal{P}\!\mathit{ic}}
\def\Picardfunctor{\mathrm{Pic}}
\def\Deformationcategory{\mathcal{D}\!\mathit{ef}}
\usepackage{cleveref}
\crefname{originaltheorem}{original theorem}{original theorems}
\Crefname{originaltheorem}{Original theorem}{Original theorems}
\makeatletter
\newlabel{bootstrap-theorem-final-bootstrap}{{04S6}{source}{}{sourcepointer-1}{}}
\newlabel{bootstrap-theorem-final-bootstrap@cref}{{[originaltheorem][0][]04S6}{[1][0][]source}{}{}{}}
\makeatother
\crefname{originalequation}{original equation}{original equations}
\Crefname{originalequation}{Original equation}{Original equations}
\makeatletter
\newlabel{equation-completion-functor}{{0880}{source}{}{sourcepointer-2}{}}
\newlabel{equation-completion-functor@cref}{{[originalequation][0][]0880}{[1][0][]source}{}{}{}}
\makeatother
\crefname{originaldefinition}{original definition}{original definitions}
\Crefname{originaldefinition}{Original definition}{Original definitions}
\makeatletter
\newlabel{formal-spaces-definition-formal-algebraic-space}{{0AIM}{source}{}{sourcepointer-3}{}}
\newlabel{formal-spaces-definition-formal-algebraic-space@cref}{{[originaldefinition][0][]0AIM}{[1][0][]source}{}{}{}}
\makeatother
\crefname{originalreference}{original reference}{original references}
\Crefname{originalreference}{Original reference}{Original references}
\makeatletter
\newlabel{lemma-algebraize-rig-etale-affine}{{0AR4}{source}{}{sourcepointer-4}{}}
\newlabel{lemma-algebraize-rig-etale-affine@cref}{{[originalreference][0][]0AR4}{[1][0][]source}{}{}{}}
\makeatother
\crefname{originalreference}{original reference}{original references}
\Crefname{originalreference}{Original reference}{Original references}
\makeatletter
\newlabel{lemma-completion-functor-fully-faithful}{{0GDR}{source}{}{sourcepointer-5}{}}
\newlabel{lemma-completion-functor-fully-faithful@cref}{{[originalreference][0][]0GDR}{[1][0][]source}{}{}{}}
\makeatother
\crefname{originalreference}{original reference}{original references}
\Crefname{originalreference}{Original reference}{Original references}
\makeatletter
\newlabel{lemma-equivalence-relation}{{024E}{source}{}{sourcepointer-6}{}}
\newlabel{lemma-equivalence-relation@cref}{{[originalreference][0][]024E}{[1][0][]source}{}{}{}}
\makeatother
\crefname{originalreference}{original reference}{original references}
\Crefname{originalreference}{Original reference}{Original references}
\makeatletter
\newlabel{lemma-smash-away-from-T}{{0AR9}{source}{}{sourcepointer-7}{}}
\newlabel{lemma-smash-away-from-T@cref}{{[originalreference][0][]0AR9}{[1][0][]source}{}{}{}}
\makeatother
\crefname{originalsection}{original section}{original sections}
\Crefname{originalsection}{Original section}{Original sections}
\makeatletter
\newlabel{schemes-example-affine-space-zero-doubled}{{01JD}{source}{}{sourcepointer-8}{}}
\newlabel{schemes-example-affine-space-zero-doubled@cref}{{[originalsection][0][]01JD}{[1][0][]source}{}{}{}}
\makeatother
\title{The completion functor is an equivalence over an affine base}
\author{The Stacks Project Authors}
\date{stacks-project:restricted.tex:local-snapshot-20261012; GNU FDL 1.2 or later}
\begin{document}
\maketitle
\noindent\textit{Editorial note.} All delivered mathematical text is the author's own, copied from the pinned Stacks Project chapter file restricted.tex; no statement or proof step was rewritten, reordered or removed.\par\medskip
\noindent\textit{Editorial note.} The counted proof is the author's complete argument: the source proves the lemma in three explicit steps (affine, separated, general), quoting the algebraization and bootstrap results for the gluing; the argument is complete\par\medskip
\noindent\textit{Editorial note.} The label lemma-dilatations-affine names the counted statement, and the proof environment immediately following it in the source is delivered verbatim as its proof.\par\medskip
\noindent\textit{Editorial note.} The counted statement is carried by the author's own lemma environment in the chapter source.\par\medskip
\noindent\textit{Allowed standard prerequisites (editorial).}\begin{itemize}
\item the completion functor from the chapter's construction and its full faithfulness lemma (counted earlier in this lane)
\item rig-etale morphisms of formal algebraic spaces and the algebraization lemma for affine rig-etale targets
\item the smash construction away from T producing formal completions of etale covers
\item etale equivalence relations and the bootstrap theorem producing quotient algebraic spaces
\item compatibility of completion with fibre products and quotient sheaves
\end{itemize}\medskip
\ifdefined\hypertarget\hypertarget{sourcepointer-1}{}\fi
\noindent\textit{Source reference.} Original theorem 04S6: sources/stacks-chow/tags.txt, source lines 6209--6209. Prerequisite pointer invoked at proof line(s) 6999; it names a standard fact of a sibling chapter that the counted argument uses, and every step of the counted proof built on it is retained verbatim in the delivered body.\par\medskip
\ifdefined\hypertarget\hypertarget{sourcepointer-2}{}\fi
\noindent\textit{Source reference.} Original equation 0880: sources/stacks-restricted/pinned-original-source.tex, source lines 6766--6766. Prerequisite pointer invoked at proof line(s) 6948; it names a standard fact of the same chapter that the counted argument uses, and every step of the counted proof built on it is retained verbatim in the delivered body.\par\medskip
\ifdefined\hypertarget\hypertarget{sourcepointer-3}{}\fi
\noindent\textit{Source reference.} Original definition 0AIM: sources/stacks-chow/tags.txt, source lines 13270--13270. Prerequisite pointer invoked at proof line(s) 6980; it names a standard fact of a sibling chapter that the counted argument uses, and every step of the counted proof built on it is retained verbatim in the delivered body.\par\medskip
\ifdefined\hypertarget\hypertarget{sourcepointer-4}{}\fi
\noindent\textit{Source reference.} Original reference 0AR4: sources/stacks-restricted/pinned-original-source.tex, source lines 6406--6406. Prerequisite pointer invoked at proof line(s) 6974; it names a standard fact of the same chapter that the counted argument uses, and every step of the counted proof built on it is retained verbatim in the delivered body.\par\medskip
\ifdefined\hypertarget\hypertarget{sourcepointer-5}{}\fi
\noindent\textit{Source reference.} Original reference 0GDR: sources/stacks-restricted/pinned-original-source.tex, source lines 6832--6832. Prerequisite pointer invoked at proof line(s) 6964; it names a standard fact of the same chapter that the counted argument uses, and every step of the counted proof built on it is retained verbatim in the delivered body.\par\medskip
\ifdefined\hypertarget\hypertarget{sourcepointer-6}{}\fi
\noindent\textit{Source reference.} Original reference 024E: sources/stacks-restricted/pinned-original-source.tex, source lines 6517--6517. Prerequisite pointer invoked at proof line(s) 6995; it names a standard fact of the same chapter that the counted argument uses, and every step of the counted proof built on it is retained verbatim in the delivered body.\par\medskip
\ifdefined\hypertarget\hypertarget{sourcepointer-7}{}\fi
\noindent\textit{Source reference.} Original reference 0AR9: sources/stacks-restricted/pinned-original-source.tex, source lines 6554--6554. Prerequisite pointer invoked at proof line(s) 6959; it names a standard fact of the same chapter that the counted argument uses, and every step of the counted proof built on it is retained verbatim in the delivered body.\par\medskip
\ifdefined\hypertarget\hypertarget{sourcepointer-8}{}\fi
\noindent\textit{Source reference.} Original section 01JD: sources/stacks-chow/tags.txt, source lines 2010--2010. Prerequisite pointer invoked at proof line(s) 6957; it names a standard fact of a sibling chapter that the counted argument uses, and every step of the counted proof built on it is retained verbatim in the delivered body.\par\medskip
\begin{lemma}
\label{lemma-dilatations-affine}
In the situation above. Assume $X$ is affine. Then the functor
(\ref{equation-completion-functor}) is an equivalence.
\end{lemma}

\begin{proof}
We already know the functor is fully faithful, see
Lemma \ref{lemma-completion-functor-fully-faithful}.
Essential surjectivity. Let $g : W \to X_{/T}$ be a morphism
of formal algebraic spaces with $W$ locally Noetherian
and $g$ rig-\'etale. We will prove $W$ is in the essential
image in a number of steps.

\medskip\noindent
Step 1: $W$ is an affine formal algebraic space. Then we can find
$U \to X$ of finite type and \'etale over $X \setminus T$ such that
$U_{/T}$ is isomorphic to $W$, see
Lemma \ref{lemma-algebraize-rig-etale-affine}.
Thus we see that $W$ is in the essential image by
Lemma \ref{lemma-smash-away-from-T}.

\medskip\noindent
Step 2: $W$ is separated. Choose $\{W_i \to W\}$ as in
Formal Spaces, Definition \ref{formal-spaces-definition-formal-algebraic-space}.
By Step 1 the formal algebraic spaces $W_i$ and $W_i \times_W W_j$
are in the essential image.
Say $W_i = (X'_i)_{/T}$ and $W_i \times_W W_j = (X'_{ij})_{/T}$.
By fully faithfulness we obtain morphisms $t_{ij} : X'_{ij} \to X'_i$
and $s_{ij} : X'_{ij} \to X'_j$ matching the projections
$W_i \times_W W_j \to W_i$ and $W_i \times_W W_j \to W_j$.
Consider the structure
$$
R = \coprod X'_{ij},\quad
V = \coprod X'_i,\quad
s = \coprod s_{ij},\quad
t = \coprod t_{ij}
$$
(We can't use the letter $U$ as it has already been used.)
Applying Lemma \ref{lemma-equivalence-relation} we find that
$(t, s) : R \to V \times_X V$ defines an \'etale equivalence relation
on $V$ over $X$. Thus we can take the quotient $X' = V/R$ and
it is an algebraic space, see
Bootstrap, Theorem \ref{bootstrap-theorem-final-bootstrap}.
Since completion commutes with fibre products and taking
quotient sheaves, we find that
$X'_{/T} \cong W$ as formal algebraic spaces over $X_{/T}$.

\medskip\noindent
Step 3: $W$ is general. Choose $\{W_i \to W\}$ as in
Formal Spaces, Definition \ref{formal-spaces-definition-formal-algebraic-space}.
The formal algebraic spaces $W_i$ and $W_i \times_W W_j$ are separated.
Hence by Step 2  the formal algebraic spaces $W_i$ and $W_i \times_W W_j$
are in the essential image. Then we argue exactly as in the previous
paragraph to see that $W$ is in the essential image as well.
This concludes the proof.
\end{proof}

\end{document}
